\documentclass[11pt,a4paper]{article}

% ── Packages ──────────────────────────────────────────────────────────────
\usepackage[margin=1in]{geometry}
\usepackage{amsmath,amssymb,amsthm}
\usepackage{graphicx}
\usepackage{booktabs}
\usepackage{hyperref}
\usepackage{xcolor}
\usepackage{listings}
\usepackage{caption}
\usepackage{subcaption}
\usepackage{enumitem}
\usepackage{microtype}
\usepackage{natbib}
\usepackage{algorithm}
\usepackage{algpseudocode}
\usepackage{tikz}
\usetikzlibrary{arrows.meta, positioning, shapes.geometric}

% ── Hyperref setup ────────────────────────────────────────────────────────
\hypersetup{
  colorlinks=true,
  linkcolor=blue!70!black,
  citecolor=green!60!black,
  urlcolor=blue!80!black
}

% ── Code listing style ────────────────────────────────────────────────────
\definecolor{codegray}{rgb}{0.5,0.5,0.5}
\definecolor{codeblue}{rgb}{0.0,0.35,0.75}
\lstset{
  basicstyle=\ttfamily\small,
  keywordstyle=\color{codeblue}\bfseries,
  commentstyle=\color{codegray}\itshape,
  numberstyle=\tiny\color{codegray},
  numbers=left,
  stepnumber=1,
  breaklines=true,
  frame=single,
  language=Python
}

% ── Theorem environments ──────────────────────────────────────────────────
\newtheorem{theorem}{Theorem}
\newtheorem{lemma}[theorem]{Lemma}
\newtheorem{corollary}[theorem]{Corollary}
\newtheorem{definition}{Definition}
\newtheorem{proposition}[theorem]{Proposition}

% ── Custom commands ───────────────────────────────────────────────────────
\newcommand{\Tcomp}{T_{\mathrm{compute}}}
\newcommand{\Tmem}{T_{\mathrm{memory}}}
\newcommand{\BW}{\mathcal{B}}
\newcommand{\FLOPS}{\mathcal{F}}
\newcommand{\eg}{\textit{e.g.}}
\newcommand{\ie}{\textit{i.e.}}

% ─────────────────────────────────────────────────────────────────────────
\begin{document}

% ── Title block ───────────────────────────────────────────────────────────
\title{%
  \textbf{The Physics of LLM Inference Cost:}\\
  \large How Hardware Constraints Dictate Pricing, Architecture,\\
  and the Economics of Large Language Model Deployment
}

\author{%
  Mohan Chinnappan\\
  \texttt{mohan.chinnappan.n@gmail.com}\\[6pt]
  \small New York City, NY
}

\date{\today}

\maketitle

% ── Abstract ──────────────────────────────────────────────────────────────
\begin{abstract}
The pricing structure of large language model (LLM) inference services —
including premium charges for low-latency modes, surcharges for long context
windows, and discounts for cached contexts — appears arbitrary from a business
perspective but is in fact entirely determined by fundamental physical laws
governing GPU hardware. In this paper we derive a formal model of LLM inference
cost grounded in hardware physics, demonstrating that runtime is the maximum of
two quantities: compute time, which scales with batch size and active parameters,
and memory time, which scales with total model weight and KV cache size. We show
that for virtually all real-world inference workloads, memory bandwidth is the
binding constraint, not arithmetic throughput. From this single equation we derive
closed-form explanations for the $6\times$ fast-mode premium, the 50\% long-context
surcharge, the 5-minute cache discount structure, the MoE architecture boundary at
the rack level, and the $100\times$ over-training phenomenon observed in frontier
models. We further analyse reversible network architectures as a mechanism for
bypassing activation memory constraints during training, and characterise the physical
limits imposed by NVLink/InfiniBand topology on model parallelism strategies.
Our analysis provides practitioners with a principled framework for predicting
inference costs, designing cost-optimal deployment configurations, and anticipating
future pricing dynamics as hardware generations evolve.

\medskip\noindent
\textbf{Keywords:} LLM inference, memory bandwidth, roofline model, KV cache,
mixture of experts, pipeline parallelism, inference cost, GPU memory hierarchy
\end{abstract}

\tableofcontents
\newpage

% ═══════════════════════════════════════════════════════════════════════════
\section{Introduction}
\label{sec:intro}
% ═══════════════════════════════════════════════════════════════════════════

The rapid deployment of large language models (LLMs) at commercial scale has
created a complex and often opaque pricing landscape. API providers charge
vastly different rates depending on response latency, context length, and
cache utilisation. A ``fast mode'' API call may cost $6\times$ more than a
standard call for identical output. A 200{,}000-token context costs
approximately 50\% more per token than an 8{,}000-token context. A cache hit
on a recently-used context may receive an 80--95\% discount.

These pricing differentials appear arbitrary from a business perspective but
are, in fact, mathematically determined by the physical constraints of modern
GPU hardware. Every pricing tier corresponds to a specific hardware
constraint: memory bandwidth limits, HBM drain times, NVLink topology, and
the thermodynamics of data movement.

\subsection{Contributions}

This paper makes the following contributions:

\begin{enumerate}[leftmargin=*]
  \item We derive the \textbf{Fundamental Inference Equation} (Section
    \ref{sec:fundamental}) from first principles, showing that inference
    runtime is determined by the maximum of compute and memory transfer times.

  \item We develop the \textbf{Cost Parabola Model} (Section \ref{sec:roofline})
    characterising the U-shaped cost-per-token curve as a function of batch
    size, and derive the optimal batch size analytically.

  \item We provide \textbf{closed-form derivations} of common pricing
    phenomena including the fast-mode premium, the context-length surcharge,
    and the cache discount tiers (Sections \ref{sec:batch}--\ref{sec:tiers}).

  \item We analyse the \textbf{physical constraints on model parallelism}
    imposed by NVLink/InfiniBand topology (Section \ref{sec:rack}), explaining
    the architectural necessity of expert parallelism boundaries.

  \item We derive the \textbf{over-training optimum} for models deployed at
    scale (Section \ref{sec:overtrain}) and analyse reversible architectures
    for training memory efficiency (Section \ref{sec:revnet}).
\end{enumerate}

\subsection{Paper Organisation}

Section \ref{sec:background} reviews relevant hardware specifications. Section
\ref{sec:fundamental} presents the fundamental inference equation. Sections
\ref{sec:roofline}--\ref{sec:tiers} derive the roofline model and memory tier
economics. Section \ref{sec:rack} analyses rack-level topology constraints.
Sections \ref{sec:moe}--\ref{sec:pipeline} cover parallelism strategies.
Sections \ref{sec:overtrain}--\ref{sec:revnet} analyse training-inference
cost coupling. Section \ref{sec:conclusion} concludes with implications for
future hardware and model design.

% ═══════════════════════════════════════════════════════════════════════════
\section{Background and Hardware Specifications}
\label{sec:background}
% ═══════════════════════════════════════════════════════════════════════════

\subsection{Modern GPU Memory Hierarchy}

Contemporary AI accelerators expose a four-level memory hierarchy, each with
dramatically different bandwidth, capacity, and cost characteristics:

\begin{table}[h]
\centering
\caption{Memory hierarchy specifications (NVIDIA H100 SXM5)}
\label{tab:memory-hierarchy}
\begin{tabular}{lrrrr}
\toprule
\textbf{Level} & \textbf{Bandwidth} & \textbf{Capacity} & \textbf{Latency} & \textbf{API Tier} \\
\midrule
HBM3 (on-chip DRAM) & 3.35 TB/s & 80 GB  & $\sim$500 cycles  & Instant response \\
DDR5 (host memory)  & 0.1 TB/s  & 1--2 TB & $\sim$100 ns     & 5-min cache \\
NVMe SSD            & 12 GB/s   & 10+ TB  & $\sim$100 $\mu$s & 1-hour cache \\
HDD / Object Store  & 0.5 GB/s  & $\infty$& $>$1 ms          & Cold start \\
\bottomrule
\end{tabular}
\end{table}

\subsection{Arithmetic Throughput}

The H100 SXM5 provides 989 TFLOPS of BF16 tensor-core throughput and
1{,}979 TFLOPS of FP8 throughput. The critical ratio characterising whether
a workload is compute-bound or memory-bound is the \textit{arithmetic
intensity}: the number of floating-point operations performed per byte of
data moved.

\begin{definition}[Arithmetic Intensity]
\label{def:arith-intensity}
For a computation performing $F$ FLOPs while moving $B$ bytes of data,
the arithmetic intensity is:
\[
  I = \frac{F}{B} \quad \text{(FLOPs/byte)}
\]
A computation with $I < I^*$ is memory-bound; otherwise compute-bound,
where $I^* = \FLOPS / \BW$ is the hardware's ridge point.
\end{definition}

For H100: $I^* = 989 \times 10^{12} / 3.35 \times 10^{12} \approx 295$ FLOPs/byte.
A single-batch transformer decode step has $I \approx 2$ FLOPs/byte — well
below the ridge point, confirming memory boundedness.

% ═══════════════════════════════════════════════════════════════════════════
\section{The Fundamental Inference Equation}
\label{sec:fundamental}
% ═══════════════════════════════════════════════════════════════════════════

\subsection{Derivation}

An LLM inference step involves two separable time components that execute
in parallel on modern hardware:

\begin{definition}[Compute Time]
\label{def:tcompute}
For a batch of $B$ sequences through a model with $P_\text{active}$ active
parameters and peak throughput $\FLOPS$:
\[
  \Tcomp = \frac{2 \cdot B \cdot P_\text{active}}{\FLOPS}
\]
The factor of 2 accounts for the multiply-accumulate (MAC) structure of
matrix multiplications.
\end{definition}

\begin{definition}[Memory Transfer Time]
\label{def:tmemory}
For model weights of total size $W_\text{bytes}$ and a KV cache of size
$C_\text{bytes}$ per sequence, with memory bandwidth $\BW$:
\[
  \Tmem = \frac{W_\text{bytes} + B \cdot C_\text{bytes}}{\BW}
\]
\end{definition}

\begin{theorem}[Fundamental Inference Equation]
\label{thm:fundamental}
The runtime of a single inference step is:
\[
  \boxed{T_\text{inference} = \max\!\left(\Tcomp,\; \Tmem\right)}
\]
The system is either compute-bound ($\Tcomp > \Tmem$) or memory-bound
($\Tmem > \Tcomp$), with the binding constraint determining total latency.
\end{theorem}

\begin{proof}
Compute and memory transfer can proceed simultaneously due to hardware
pipelining. The slower operation determines completion time; hence the
maximum. \qed
\end{proof}

\subsection{Regime Analysis}

Setting $\Tcomp = \Tmem$ and solving for $B$ gives the critical batch size:

\begin{proposition}[Critical Batch Size]
\label{prop:critical-batch}
The batch size at which compute and memory become co-limiting is:
\[
  B^* = \frac{\FLOPS \cdot (W_\text{bytes} + B^* \cdot C_\text{bytes})}
              {2 \cdot P_\text{active} \cdot \BW}
\]
For $C_\text{bytes} \ll W_\text{bytes}/B^*$ (short contexts):
\[
  B^* \approx \frac{\FLOPS \cdot W_\text{bytes}}
                   {2 \cdot P_\text{active} \cdot \BW}
             = \frac{I^* \cdot W_\text{bytes}}{2 \cdot P_\text{active}}
\]
\end{proposition}

\textbf{Numerical example (H100, Llama-3 70B FP16):}
\[
  B^* \approx \frac{295 \times 140 \times 10^9}{2 \times 70 \times 10^9}
            = \frac{41{,}300}{140} \approx 295 \text{ sequences}
\]

This matches the empirically observed optimal batch size of approximately
300 concurrent sequences for dense 70B models.

% ═══════════════════════════════════════════════════════════════════════════
\section{The Roofline Model and Cost Parabola}
\label{sec:roofline}
% ═══════════════════════════════════════════════════════════════════════════

\subsection{Cost-Per-Token as a Function of Batch Size}

Let $\lambda$ be the cost per GPU-second. The cost per output token is:
\[
  \text{cost}(B) = \frac{\lambda \cdot T_\text{inference}(B)}{B}
                = \frac{\lambda}{B} \cdot \max\!\left(\Tcomp(B),\, \Tmem(B)\right)
\]

\begin{theorem}[Cost Parabola]
\label{thm:parabola}
For a fixed context length, cost-per-token as a function of batch size is:
\[
  \text{cost}(B) = \lambda \cdot
  \begin{cases}
    \dfrac{W_\text{bytes}}{B \cdot \BW} + \dfrac{C_\text{bytes}}{\BW}
      & \text{if } B < B^* \quad (\text{memory-bound}) \\[10pt]
    \dfrac{2 P_\text{active}}{\FLOPS}
      & \text{if } B \geq B^* \quad (\text{compute-bound})
  \end{cases}
\]
The minimum is achieved at $B = B^*$, and the curve is convex, forming a
characteristic ``parabolic'' shape.
\end{theorem}

\textbf{Implication — the Fast-Mode Premium:}
A provider offering batch size $B=1$ (immediate response) incurs a cost
proportional to $W_\text{bytes}/\BW$ entirely borne by one token.
At optimal batch $B^*$, this cost is amortised over $B^*$ tokens.
The premium ratio is approximately:
\[
  \frac{\text{cost}(1)}{\text{cost}(B^*)}
  = \frac{W_\text{bytes}/\BW}
         {W_\text{bytes}/(B^* \cdot \BW)}
  = B^* \approx 300
\]
In practice, providers observe $\approx 6\times$ premiums because batch
filling is imperfect and the model also incurs per-token compute costs.

% ═══════════════════════════════════════════════════════════════════════════
\section{The 20ms Batching Schedule}
\label{sec:batch}
% ═══════════════════════════════════════════════════════════════════════════

\subsection{HBM Drain Time as the Fundamental Clock}

The HBM drain time is the time required to stream all model weights once:
\[
  t_\text{drain} = \frac{W_\text{bytes}}{\BW}
                 = \frac{140 \times 10^9 \text{ bytes}}{3.35 \times 10^{12} \text{ bytes/s}}
                 \approx 41.8 \text{ ms}
\]

An inference server structured for maximum throughput schedules one batch
per drain cycle, departing every $t_\text{drain} \approx 42$ ms. A request
that arrives just after a batch departs waits at most $2 \times t_\text{drain}
\approx 84$ ms in the worst case.

\subsection{Optimal Batch Formula}

Let $\sigma$ denote the MoE sparsity factor (fraction of parameters activated
per token; $\sigma=1$ for dense models). The economically optimal batch size is:

\begin{equation}
  B_\text{opt} \approx 300 \cdot \sigma
  \label{eq:optimal-batch}
\end{equation}

For a dense model ($\sigma=1$): $B_\text{opt} \approx 300$ sequences.
For MoE with top-2 of 64 experts ($\sigma = 2/64 \approx 0.031$):
$B_\text{opt} \approx 9.4$, but with 64× fewer active parameters,
the per-token cost is equivalent to a dense model batch of 300.

% ═══════════════════════════════════════════════════════════════════════════
\section{The Context Length Penalty}
\label{sec:context}
% ═══════════════════════════════════════════════════════════════════════════

\subsection{KV Cache Memory Footprint}

For a transformer with $L$ layers, $H_\text{kv}$ key-value heads of
dimension $d_h$, processing a context of $S$ tokens in FP16:

\begin{equation}
  C_\text{bytes}(S) = 2 \cdot L \cdot H_\text{kv} \cdot d_h \cdot S \cdot 2
  \label{eq:kv-cache}
\end{equation}

The factor of 2 at the front accounts for keys and values; the trailing 2
is bytes per FP16 element.

\textbf{Llama-3 70B:} $L=80$, $H_\text{kv}=8$, $d_h=128$:
\[
  C_\text{bytes}(S) = 2 \times 80 \times 8 \times 128 \times S \times 2
                    = 327{,}680 \cdot S \text{ bytes} \approx 0.328 \text{ MB/token}
\]

\subsection{Memory Time with Long Context}

Substituting into $\Tmem$:
\[
  \Tmem(B, S) = \frac{W_\text{bytes} + B \cdot C_\text{bytes}(S)}{\BW}
\]

The ratio of KV cache to model weights:
\[
  \rho(B, S) = \frac{B \cdot C_\text{bytes}(S)}{W_\text{bytes}}
             = \frac{B \cdot 0.328 \text{ MB} \cdot S}{140{,}000 \text{ MB}}
\]

At $S = 10^6$ tokens, $B=1$: $\rho \approx 2.34$ — the KV cache is
$2.34\times$ larger than the model weights, completely dominating inference time.

\begin{corollary}[Context Length Surcharge]
The cost increase from context length $S_0$ to $S_1$ at fixed $B$ is:
\[
  \frac{\text{cost}(S_1)}{\text{cost}(S_0)}
  \approx \frac{W_\text{bytes} + B \cdot C_\text{bytes}(S_1)}
               {W_\text{bytes} + B \cdot C_\text{bytes}(S_0)}
\]
For $S_0 = 8{,}192$, $S_1 = 200{,}000$, $B=300$: the ratio $\approx 1.49$,
predicting approximately the observed 50\% surcharge.
\end{corollary}

% ═══════════════════════════════════════════════════════════════════════════
\section{Memory Tier Economics and Cache Pricing}
\label{sec:tiers}
% ═══════════════════════════════════════════════════════════════════════════

\subsection{Tiered Storage Model}

Let $\mathcal{T} = \{$HBM, DDR, Flash$\}$ denote the storage tiers.
Each tier $i$ has bandwidth $\BW_i$, holding cost $\lambda_i$ per GB per hour,
and transfer time $t_i = C/\BW_i$ to load a KV cache of size $C$ bytes.

\begin{table}[h]
\centering
\caption{Memory tier economics}
\label{tab:tiers}
\begin{tabular}{lrrrr}
\toprule
\textbf{Tier} & \textbf{Bandwidth} & \textbf{Drain Time (100GB)} & \textbf{Cost (\$/GB-hr)} & \textbf{API} \\
\midrule
HBM   & 3.35 TB/s & 30 ms    & $\approx$\$30  & Instant \\
DDR   & 100 GB/s  & 1 s      & $\approx$\$0.03 & 5-min discount \\
Flash & 12 GB/s   & 8 s      & $\approx$\$0.003 & 1-hr discount \\
\bottomrule
\end{tabular}
\end{table}

\subsection{Optimal Eviction Policy}

A provider maximises profit by evicting a KV cache to cheaper storage when
the expected cost of keeping it in HBM exceeds the expected reload cost.
Let $p(t)$ be the probability that the user makes another request within
time $t$. The eviction threshold is:

\[
  t^*_i = \frac{C \cdot (t_\text{reload,i} / t_\text{drain})}
               {(\lambda_\text{HBM} - \lambda_i) \cdot C / 3600}
\]

For HBM $\to$ DDR with $t_\text{reload} = 1$ s and
$\Delta\lambda \approx \$30$/GB-hr:
\[
  t^* \approx \frac{1/42 \times 10^{-3}}{(30 - 0.03)/3600} \approx 290 \text{ s}
\]
This aligns precisely with the 5-minute (300 s) cache tier offered commercially.

% ═══════════════════════════════════════════════════════════════════════════
\section{The Physical Rack Limit and Network Topology}
\label{sec:rack}
% ═══════════════════════════════════════════════════════════════════════════

\subsection{Intra-Rack vs. Inter-Rack Bandwidth}

Within a DGX H100 node, NVLink 4.0 provides 900 GB/s all-to-all bandwidth
via NVSwitch. Crossing rack boundaries requires InfiniBand NDR at 400 Gbps
(50 GB/s per port). The bandwidth ratio is:

\[
  \frac{\BW_\text{NVLink}}{\BW_\text{IB}} \approx \frac{900}{50} = 18\times
\]

\subsection{Physical Constraints on NVLink Scaling}

The physical reason NVLink cannot scale beyond a single rack is fundamental:
NVLink uses passive copper cables, each with a minimum bend radius of
$\sim$15 mm. In an 8-GPU chassis, the required cable density already consumes
a significant fraction of the available cross-sectional area. Adding more
GPUs to a single NVLink domain requires quadratically more cables, bounded
by finite physical volume, cable weight limits, and thermal constraints.

\begin{proposition}[Rack Topology Constraint]
Any all-to-all communication pattern (\eg, MoE expert routing) achieves
maximum efficiency when confined within a single NVLink domain. The effective
bandwidth of an all-to-all operation spanning $r$ racks is:
\[
  \BW_\text{eff}(r) = \frac{\BW_\text{NVLink}}{1 + (r-1) \cdot k}
\]
where $k \approx 18$ is the NVLink/InfiniBand bandwidth ratio. For $r=2$,
$\BW_\text{eff} \approx \BW_\text{NVLink}/19$, an approximate $19\times$ penalty.
\end{proposition}

% ═══════════════════════════════════════════════════════════════════════════
\section{Mixture of Experts as a Hardware-Native Architecture}
\label{sec:moe}
% ═══════════════════════════════════════════════════════════════════════════

\subsection{MoE as a Rack Mirror}

In an MoE model with $E$ experts and top-$k$ routing, each token is routed
to $k$ of $E$ experts, which may reside on any GPU in the group. The
resulting communication pattern is all-to-all: any GPU may need to send
token activations to any other GPU. This is precisely the NVLink-native
pattern.

\begin{theorem}[MoE Rack Boundary]
\label{thm:moe-rack}
For a MoE layer with $E$ experts distributed across $G$ GPUs in a NVLink
domain and $N$ tokens, the all-to-all communication volume per layer is:
\[
  V_\text{all2all} = N \cdot d_\text{model} \cdot k \cdot 2 \cdot \text{dtype\_bytes}
\]
This is minimised (and bandwidth is maximised) when the entire expert set
fits within a single NVLink domain, with the physical rack as the natural
boundary.
\end{theorem}

\subsection{Memory Efficiency of MoE}

Let $P_\text{total}$ be total parameters and $P_\text{active} = (k/E) \cdot P_\text{total}$
be active parameters per token. The memory bandwidth saving relative to a
dense model of equal total size:
\[
  \text{BW Saving} = \frac{P_\text{total} - P_\text{active}}{P_\text{total}}
                   = 1 - \frac{k}{E}
\]
For top-2 of 64 experts: saving $= 1 - 2/64 = 96.9\%$. The model effectively
requires $3.1\times$ less memory bandwidth per token than a dense model of
equal total parameter count.

% ═══════════════════════════════════════════════════════════════════════════
\section{Pipeline Parallelism and the Bubble Cost}
\label{sec:pipeline}
% ═══════════════════════════════════════════════════════════════════════════

\subsection{The Pipeline Bubble}

When a model with $N_L$ layers is distributed across $p$ pipeline stages
(racks), each stage contains $N_L/p$ layers. Processing a batch requires
micro-batching to fill the pipeline. With $m$ micro-batches and $p$ stages,
the pipeline bubble fraction is:

\begin{equation}
  \phi = \frac{p - 1}{m + p - 1}
  \label{eq:bubble}
\end{equation}

\textbf{GPU utilisation:} $U_\text{pipe} = 1 - \phi = \frac{m}{m + p - 1}$

\begin{corollary}[Minimum Micro-Batch Requirement]
To achieve GPU utilisation $U \geq U_\text{min}$, the minimum number of
micro-batches required is:
\[
  m \geq \frac{U_\text{min} \cdot (p-1)}{1 - U_\text{min}}
\]
For $U_\text{min} = 0.9$ and $p = 4$ racks: $m \geq 27$ micro-batches.
\end{corollary}

% ═══════════════════════════════════════════════════════════════════════════
\section{The Over-Training Optimum}
\label{sec:overtrain}
% ═══════════════════════════════════════════════════════════════════════════

\subsection{Compute Balance Across the Model Lifecycle}

Let $C_\text{train}$, $C_\text{align}$, and $C_\text{infer}$ denote the
total compute consumed during pre-training, alignment (RLHF), and inference
deployment respectively. The total ecosystem cost is minimised when:

\begin{equation}
  C_\text{train} = C_\text{align} = C_\text{infer}
  \label{eq:compute-balance}
\end{equation}

\subsection{Deriving the Over-Training Factor}

The Chinchilla-optimal \citep{hoffmann2022training} token count for a
model with $P$ parameters is $D^* \approx 20P$. However, a provider
expecting $V$ daily inference tokens over $T$ years of deployment will
over-train by a factor:

\begin{equation}
  \alpha = \frac{D_\text{actual}}{D^*}
         \approx \frac{V \cdot T \cdot C_\text{per-token-inference}}
                      {C_\text{train} / D^*}
  \label{eq:overtrain-factor}
\end{equation}

For $V = 100 \times 10^6$ daily tokens, $T = 3$ years:
\[
  \alpha \approx \frac{100 \times 10^6 \times 365 \times 3 \times 2P}
                      {6P \cdot D^*}
       = \frac{2.19 \times 10^{11}}{6 \times 20} \approx 1.8 \times 10^9
\]

In practice the over-training ratio is bounded by diminishing returns to
additional training data. Empirically, frontier labs train at
$\alpha \approx 100$ past the Chinchilla optimum, corresponding to training
on $\sim 200$ trillion tokens for a 70B model.

\textbf{Economic interpretation:} Over-training compresses more world
knowledge into the same parameter count, reducing the effective inference
cost per token. The upfront compute over-spend is amortised over the
expected inference volume.

% ═══════════════════════════════════════════════════════════════════════════
\section{Reversible Networks and the Activation Memory Wall}
\label{sec:revnet}
% ═══════════════════════════════════════════════════════════════════════════

\subsection{Standard Backpropagation Memory Cost}

Standard backpropagation must store all intermediate activations during the
forward pass for use in the backward pass. For a model with $L$ layers,
batch size $B$, sequence length $S$, and hidden dimension $d$:

\[
  M_\text{activations} = L \cdot B \cdot S \cdot d \cdot \text{dtype\_bytes}
\]

For GPT-4 scale ($L=120$, $d=12{,}288$, $S=8{,}192$, $B=1$):
$M_\text{activations} \approx 120 \times 1 \times 8{,}192 \times 12{,}288 \times 2
\approx 24$ TB — clearly impossible without gradient checkpointing.

\subsection{Reversible Networks}

A \textit{reversible layer} \citep{gomez2017reversible} is a bijective
function $f: \mathbb{R}^{2d} \to \mathbb{R}^{2d}$ that can be inverted
without storing intermediate values. The standard Feistel construction:

\begin{align}
  (y_1, y_2) &= f(x_1, x_2) \\
  y_1 &= x_1 + \mathcal{F}(x_2) \\
  y_2 &= x_2 + \mathcal{G}(y_1)
\end{align}

Inversion:
\begin{align}
  x_2 &= y_2 - \mathcal{G}(y_1) \\
  x_1 &= y_1 - \mathcal{F}(x_2)
\end{align}

\begin{theorem}[Activation Memory of Reversible Networks]
A reversible network with $L$ layers requires $O(1)$ activation memory
during backpropagation, independent of $L$, by recomputing all intermediate
activations from the final output during the backward pass:
\[
  M_\text{rev} = O(B \cdot S \cdot d) \quad \text{(constant in } L\text{)}
\]
\end{theorem}

The trade-off is a $\sim 33\%$ increase in backward pass compute, which is
typically acceptable when HBM memory is the binding constraint.

% ═══════════════════════════════════════════════════════════════════════════
\section{Practical Cost Model}
\label{sec:practical}
% ═══════════════════════════════════════════════════════════════════════════

\subsection{Unified Cost Formula}

Combining all components, the cost per output token for a dense model is:

\begin{equation}
  \text{cost}_\text{token}(B, S) = \frac{\lambda_\text{GPU}}{B} \cdot
  \max\!\left(
    \frac{2 P_\text{active}}{\FLOPS},\;
    \frac{W_\text{bytes} + B \cdot C(S)}{\BW}
  \right)
  \label{eq:unified-cost}
\end{equation}

where $C(S) = 2 L H_\text{kv} d_h S \cdot \text{dtype\_bytes}$
(Equation~\ref{eq:kv-cache}).

\subsection{Numerical Predictions vs. Observed Pricing}

\begin{table}[h]
\centering
\caption{Model predictions vs. observed commercial pricing (illustrative)}
\label{tab:predictions}
\begin{tabular}{lcc}
\toprule
\textbf{Pricing Feature}        & \textbf{Predicted Ratio} & \textbf{Observed Range} \\
\midrule
Fast mode vs. standard          & $B^* / 1 \approx 300\times$  & $4$--$8\times$  \\
                                & (amortisation ratio)      & (partial batching)\\
200k vs. 8k context surcharge   & $1.49\times$ (at $B=300$) & $1.3$--$1.7\times$ \\
5-min cache discount            & $\BW_\text{HBM}/\BW_\text{DDR} \approx 33\times$ & $5$--$20\times$ \\
MoE vs. dense cost ratio        & $k/E \approx 0.031$       & $0.05$--$0.15\times$ \\
\bottomrule
\end{tabular}
\end{table}

The discrepancy between predicted and observed fast-mode ratios reflects
two factors: (1) providers partially amortise costs even in fast mode via
continuous batching, and (2) compute costs are non-negligible at very short
sequences.

% ═══════════════════════════════════════════════════════════════════════════
\section{Related Work}
\label{sec:related}
% ═══════════════════════════════════════════════════════════════════════════

The roofline model \citep{williams2009roofline} provides the foundational
framework for characterising compute- vs. memory-bound workloads. Its
application to deep learning has been explored extensively for training
\citep{jain2019gist}, but inference has received less formal treatment.

Chinchilla scaling laws \citep{hoffmann2022training} established optimal
training compute allocation; our work extends this to the inference regime.
Orca \citep{yu2022orca} introduced continuous batching for LLM serving,
which corresponds to the practical realisation of our batching schedule.
vLLM \citep{kwon2023efficient} introduced paged attention for KV cache
management, directly addressing the context-length memory constraints we
formalise.

FlashAttention \citep{dao2022flashattention} reduces the HBM memory traffic
of attention from $O(N^2)$ to $O(N)$, consistent with our analysis that
memory bandwidth is the binding constraint. Speculative decoding
\citep{leviathan2023fast} implicitly exploits the compute-memory balance
point, using a small draft model to amortise the weight-fetch cost of a
large model.

% ═══════════════════════════════════════════════════════════════════════════
\section{Conclusion}
\label{sec:conclusion}
% ═══════════════════════════════════════════════════════════════════════════

We have demonstrated that the full range of LLM inference pricing phenomena —
fast-mode premiums, context-length surcharges, cache discount tiers, and
the architectural necessity of MoE expert parallelism within rack boundaries —
can be derived from a single physical equation:

\[
  T_\text{inference} = \max\!\left(
    \frac{2 B P_\text{active}}{\FLOPS},\;
    \frac{W_\text{bytes} + B \cdot C_\text{bytes}(S)}{\BW}
  \right)
\]

The binding constraint for virtually all real-world LLM inference is memory
bandwidth, not arithmetic throughput. This has three primary implications:

\begin{enumerate}[leftmargin=*]
  \item \textbf{For system designers:} optimise for memory bandwidth
    utilisation, not FLOPs; batch aggressively; use quantisation to reduce
    weight size; tier the KV cache to cheaper storage for idle sessions.

  \item \textbf{For model architects:} design model structures that confine
    all-to-all communication (MoE, tensor parallelism) within NVLink domains;
    use pipeline parallelism only across rack boundaries; prefer reversible
    architectures when activation memory constrains training scale.

  \item \textbf{For economists:} the memory wall, not compute shortage, is
    the primary pricing lever. Hardware generations that improve HBM
    bandwidth (H100 $\to$ H200 $\to$ B200) will reduce inference costs
    more than equivalent FLOP increases.

\end{enumerate}

The future of AI infrastructure efficiency lies not in faster arithmetic
but in \textit{hiding, compressing, and bypassing the memory wall}.

% ── Bibliography ──────────────────────────────────────────────────────────
\bibliographystyle{plainnat}
\begin{thebibliography}{99}

\bibitem[Dao et al.(2022)]{dao2022flashattention}
Dao, T., Fu, D.Y., Ermon, S., Rudra, A., and R{\'e}, C. (2022).
\textit{FlashAttention: Fast and Memory-Efficient Exact Attention with
IO-Awareness.} NeurIPS 2022.

\bibitem[Gomez et al.(2017)]{gomez2017reversible}
Gomez, A.N., Ren, M., Urtasun, R., and Grosse, R.B. (2017).
\textit{The Reversible Residual Network: Backpropagation Without Storing
Activations.} NeurIPS 2017.

\bibitem[Hoffmann et al.(2022)]{hoffmann2022training}
Hoffmann, J., Borgeaud, S., Mensch, A., \textit{et al.} (2022).
\textit{Training Compute-Optimal Large Language Models.}
arXiv:2203.15556.

\bibitem[Jain et al.(2019)]{jain2019gist}
Jain, A., Oehmke, P., and Shalf, J. (2019).
\textit{Gist: Efficient Data Encoding for Deep Neural Network Training.}
ISCA 2018.

\bibitem[Kwon et al.(2023)]{kwon2023efficient}
Kwon, W., Li, Z., Zhuang, S., \textit{et al.} (2023).
\textit{Efficient Memory Management for Large Language Model Serving with
PagedAttention.} SOSP 2023.

\bibitem[Leviathan et al.(2023)]{leviathan2023fast}
Leviathan, Y., Kalman, M., and Matias, Y. (2023).
\textit{Fast Inference from Transformers via Speculative Decoding.}
ICML 2023.

\bibitem[Williams et al.(2009)]{williams2009roofline}
Williams, S., Waterman, A., and Patterson, D. (2009).
\textit{Roofline: An Insightful Visual Performance Model for Multicore
Architectures.} Communications of the ACM, 52(4):65--76.

\bibitem[Yu et al.(2022)]{yu2022orca}
Yu, G., Kim, J., Shin, H., \textit{et al.} (2022).
\textit{Orca: A Distributed Serving System for Transformer-Based Generative
Models.} OSDI 2022.

\end{thebibliography}

% ── Appendix ──────────────────────────────────────────────────────────────
\appendix

\section{Derivation of the Optimal Batch Size Formula}
\label{app:batch-derivation}

Setting $\Tcomp(B^*) = \Tmem(B^*)$ and solving for $B^*$:
\begin{align}
  \frac{2 B^* P_\text{active}}{\FLOPS}
  &= \frac{W_\text{bytes} + B^* C_\text{bytes}}{\BW} \\
  2 B^* P_\text{active} \BW
  &= \FLOPS (W_\text{bytes} + B^* C_\text{bytes}) \\
  B^*(2 P_\text{active} \BW - \FLOPS C_\text{bytes})
  &= \FLOPS W_\text{bytes} \\
  B^*
  &= \frac{\FLOPS \cdot W_\text{bytes}}
          {2 P_\text{active} \BW - \FLOPS \cdot C_\text{bytes}}
\end{align}

When $\FLOPS \cdot C_\text{bytes} \ll 2 P_\text{active} \BW$ (short contexts):
\[
  B^* \approx \frac{\FLOPS \cdot W_\text{bytes}}{2 P_\text{active} \BW}
            = \frac{I^* \cdot W_\text{bytes}}{2 P_\text{active}}
\]

\section{Python Reference Implementation}
\label{app:code}

\begin{lstlisting}[caption={Roofline cost model (Python)}, label=lst:roofline]
import numpy as np

# H100 SXM5 hardware constants
FLOPS_PER_SEC = 989e12     # 989 TFLOPS FP16 peak
MEM_BW_BYTES  = 3.35e12    # 3.35 TB/s HBM3
GPU_COST_PER_HR = 3.0      # USD/hr (approximate spot price)

def inference_cost(
    n_params: float,       # total parameters
    batch_size: int,       # concurrent sequences
    context_len: int,      # tokens of context per sequence
    n_layers: int = 80,    # number of transformer layers
    n_kv_heads: int = 8,   # KV attention heads
    d_head: int = 128,     # head dimension
    dtype_bytes: int = 2,  # 2 for FP16, 1 for INT8
) -> dict:
    """
    Compute inference runtime and cost using the Fundamental Inference Equation.
    Returns: dict with T_compute, T_memory, runtime, cost_per_token.
    """
    # Weight size in bytes
    W_bytes = n_params * dtype_bytes

    # KV cache per sequence (K+V for all layers)
    kv_per_seq = 2 * n_layers * n_kv_heads * d_head * context_len * dtype_bytes

    # Compute time: 2 FLOPs per weight per sequence (matmul)
    T_compute = (2 * batch_size * n_params) / FLOPS_PER_SEC

    # Memory time: fetch weights + KV cache for all sequences
    T_memory  = (W_bytes + batch_size * kv_per_seq) / MEM_BW_BYTES

    # Fundamental equation: runtime = max of the two
    runtime   = max(T_compute, T_memory)
    bottleneck = "compute" if T_compute > T_memory else "memory"

    # Cost per output token (GPU cost / tokens per second)
    cost_per_token = (GPU_COST_PER_HR / 3600) * runtime / batch_size

    return {
        "T_compute_ms"  : T_compute * 1000,
        "T_memory_ms"   : T_memory  * 1000,
        "runtime_ms"    : runtime   * 1000,
        "bottleneck"    : bottleneck,
        "cost_per_token": cost_per_token,
    }

# Example: Llama-3 70B at various batch sizes
for B in [1, 10, 100, 295, 1000, 5000]:
    r = inference_cost(70e9, B, 2048)
    print(f"B={B:5d}: {r['runtime_ms']:6.1f}ms  "
          f"{r['bottleneck']:8s}  "
          f"${r['cost_per_token']*1e6:.2f}/Mtok")
\end{lstlisting}

\end{document}
